\section{习题 2.3：直线的方程（13 题）}\label{sec:23}
\tip{\textbf{核心三种表示：}已知一点 $P_0$ 与方向 $u=(a,b,c)\ne0$，有点向式 $P=P_0+tu$；当方向分量均非零可用对称式 $\frac{x-x_0}{a}=\frac{y-y_0}{b}=\frac{z-z_0}{c}$；也可用两张不平行平面的公共解表示。方向分量为 $0$ 时应使用相应坐标恒等式，不能把 $0$ 当作除数。}
\begin{center}
\begin{tikzpicture}[scale=.8,>=Stealth]
\draw[blue!50,fill=blue!8] (.5,.25)--(5.6,.7)--(6.8,2.9)--(1.7,2.4)--cycle;
\draw[teal!60,fill=teal!12,opacity=.6] (2.5,-.3)--(4.1,.2)--(4.9,3.2)--(3.3,2.6)--cycle;
\draw[very thick,red,->] (3.2,-.5)--(5.2,3.6) node[above]{$l=\pi_1\cap\pi_2$};
\node at (1.0,.55) {$\pi_1$};\node at (3.8,2.4) {$\pi_2$};
\node at (9.2,1.7) {$u=n_1\times n_2$};
\end{tikzpicture}
\end{center}

\Q{1}{根据一点和方向写出直线方程}
(1) 过 $M(2,-1,3)$，方向 $u=(1,0,3)$，参数式是
\[x=2+t,\qquad y=-1,\qquad z=3+3t.\]
对称式可写作 $\frac{x-2}{1}=\frac{z-3}{3},\ y=-1$，形式化常写为含零分母的三段式，实际应按此理解。

(2) 过 $M(2,-1,3),N(2,1,2)$，方向 $N-M=(0,2,-1)$，故
\[x=2,\qquad y=-1+2t,\qquad z=3-t.\]
\ans{与书后给出的两组直线方程一致。}

\Q{2}{两平面交线化为点向式}
\idea{若直线给为 $n_1\cdot X+d_1=0,n_2\cdot X+d_2=0$，先找一个公共点 $P_0$，再令 $u=n_1\times n_2$。}
\begin{enumerate}[label=(\arabic*)]
\item $x+y-z-1=0,\ y+2z=0$：取 $z=0$ 得 $P_0=(1,0,0)$；
\[n_1=(1,1,-1),\ n_2=(0,1,2),\ u=n_1\times n_2=(3,-2,1).\]
\ans{$\displaystyle\frac{x-1}{3}=\frac{y}{-2}=\frac{z}{1}$。}
\item $3x-y-z-1=0,\ 4y+3z+1=0$：联立取 $z=1$ 得 $y=-1,x=1/3$；方向是
\[n_1\times n_2=(1,-9,12).\]
\ans{$\displaystyle\frac{x-1/3}{1}=\frac{y+1}{-9}=\frac{z-1}{12}$。}
\end{enumerate}

\Q{3}{直线与平面的平行、相交及包含}
统一判别：若直线 $P=P_0+tu$，平面为 $n\cdot X+d=0$，则
\[
n\cdot(P_0+tu)+d=(n\cdot u)t+(n\cdot P_0+d)=0.
\]
\textbf{若 $n\cdot u\ne0$}，交点唯一；若 $n\cdot u=0$，再分 $n\cdot P_0+d=0$（直线在平面上）与不为零（严格平行）。对题设六小问依次代入即可：
\begin{enumerate}[label=(\arabic*)]
\item $u=(-1,1,1),n=(1,-2,3)$，内积为 $-1-2+3=0$，代 $P_0=(0,1,1)$ 得平面值 $6\ne0$；故严格平行。
\item $u=(1,1,2),n=(1,1,-1)$，内积 $0$；代 $P_0=(1,1,8)$ 得 $1+1-8+6=0$，故直线在平面内。
\item $u=(2,-3,-1),n=(3,2,1)$，内积 $6-6-1=-1$，故唯一相交；代参数求得交点 $(-1,2,-1)$。
\item 两平面交线的方向向量为 $n_1\times n_2$；代入给定第三平面，可得交点 $(12,-5,-1)$。
\item 将 $5x+2y-7=0$ 与 $x-y+z-6=0$ 作直线条件，检查第三平面 $4x+3y-z-1=0$，可发现直线上的每一点都满足它，故直线在平面内。
\item 对 $3x-3y+z+4=0$ 与 $x+y+z-2=0$ 的交线，平面 $3y+z-5=0$ 恰由前两张平面线性组合得到，故直线也在该平面内。
\end{enumerate}
\ans{结论依次为：平行、包含、交于 $(-1,2,-1)$、交于 $(12,-5,-1)$、包含、包含；与原书一致。}

\Q{4}{过给定直线再添加条件确定平面}
\idea{“过直线”意味着新平面方程可写成原直线两平面方程的非零线性组合：$\lambda F_1+\mu F_2=0$。然后用外加点、平行条件确定 $\lambda:\mu$。}
\begin{enumerate}[label=(\arabic*)]
\item 过 $\frac{x-1}{2}=\frac y1=\frac z{-1}$，平行 $u=(2,1,-2)$。线上的方向 $v=(2,1,-1)$，故法向量与 $v\times u=(-1,2,0)$ 同向；用点 $(1,0,0)$ 得 $-x+2y+1=0$。
\item 过 $\frac{x+2}3=\frac{y-1}0=\frac z1$ 和原点。线上取 $P=(-2,1,0)$、方向 $u=(3,0,1)$，法向量 $P\times u=(1,2,-3)$；\ans{$x+2y-3z=0$。}
\item 过两平面交线 $2x+3y+z-1=0$ 与 $x+2y-z+2=0$，平行 $u=(1,1,-1)$：把 $\lambda F_1+\mu F_2$ 的法向量与 $u$ 点乘置零并解 $\lambda:\mu$；\ans{$x+y+2z-3=0$。}
\item 过同一交线及点 $(2,2,-1)$；代点确定方程系数，\ans{$10x+11y+17z-25=0$。}
\item 过同一交线，与 $y,z$ 轴的截距数值相等：先设 $\lambda F_1+\mu F_2=0$，用截距条件 $B=C$（两个截距均存在），\ans{$4x+5y+5z-7=0$。}
\item 平行于 $6x-2y+3z+15=0$，且点 $(0,-2,-1)$ 到新旧平面的距离相等。设新平面 $6x-2y+3z+D=0$；给定点在线性式 $6x-2y+3z=1$，故 $|1+D|=|1+15|=16$，取与旧平面不同的 $D=-17$，得到\ans{$6x-2y+3z-17=0$。}
\end{enumerate}

\Q{5}{过直线且过原点或平行某坐标轴的平面}
给定直线 $F_1=4x-y+3z-1=0$、$F_2=x+5y-z+2=0$。过它的平面族为
\[\lambda F_1+\mu F_2=0,\quad (\lambda,\mu)\ne(0,0).\]
(1) 过原点：$-\lambda+2\mu=0$，取 $\lambda=2,\mu=1$ 得
\ans{$9x+3y+5z=0$。}

(2) 平行 $y$ 轴：$-\lambda+5\mu=0$，取 $\lambda=5,\mu=1$ 得
\ans{$21x+14z-3=0$。}

\Q{6}{两直线的位置关系}
\idea{先用方向外积判平行与否，再检查两直线是否共面。互不平行且共面即相交；互不平行且不共面即异面。}
\begin{enumerate}[label=(\arabic*)]
\item 方向 $u_1=(3,9,1)$ 与 $u_2=(-1,2,3)$ 不成比例。取两线上点代入 $\det(P_2-P_1,u_1,u_2)$ 为零，故\textbf{相交}。
\item 两平面交线的方向由 $n_1\times n_2$ 得出，与给定直线方向 $(1,-1,2)$ 组合检查；混合积不为零，故\textbf{异面}。
\item 两组平面各自交线得到方向相同，且包含相同一点，故\textbf{重合}。
\end{enumerate}
\ans{(1) 相交，(2) 异面，(3) 重合；与原书一致。}

\Q{7}{确定直线的点向式参数}
已知目标直线是
\[3x-y+z-2=0,\qquad x-2z+1=0\]
的公共解。第二式给 $x=2z-1$，代入第一式得 $y=7z-5$；以 $z$ 为参数，
\[(x,y,z)=(-1,-5,0)+z(2,7,1).\]
将 $z=0$ 的点和方向代回目标点向式 $\frac{x+3}{a}=\frac{y-s}{b}=\frac{z-t}{1}$，得到
\ans{$a=2,b=7,s=-12,t=-1$，与原书一致。}

\Q{8}{过点、交线、平行条件综合求直线}
\begin{enumerate}[label=(\arabic*)]
\item 过 $P=(0,1,-1)$，与平面 $x-3y+z-2=0$ 平行，同时与由 $3x-2y+2z+3=0$、$2x+y+z+1=0$ 决定的直线共面。

设目标方向 $v$，先由法向量 $n=(1,-3,1)$ 知 $n\cdot v=0$。再取已知直线方向 $u=n_1\times n_2$ 和线上一点 $Q$，共面条件为 $\det(Q-P,u,v)=0$。解得目标方向（由两条齐次约束的零空间确定），写成两平面的交线更不易丢解：
\ans{$\begin{cases}x-3y+z+4=0,\\5x-y+3z+4=0.\end{cases}$}

\item 与 $u=(8,7,1)$ 平行，且同时穿过题中两条给定直线。设第一直线上点 $P(t)=(-13+2t,5+3t,t)$，第二直线上点 $Q(s)=(10+5s,-7+4s,s)$。要求 $Q(s)-P(t)=\lambda(8,7,1)$，联立三分量方程解得 $s,t,\lambda$，写出直线一般式：
\ans{$\begin{cases}2x-3y+5z+41=0,\\x-y-z-17=0.\end{cases}$}

\item 过 $P=(0,1,-1)$，与 $\{2x-y-5=0,3x-2z+7=0\}$、$\{x+5y-10=0,y+z-3=0\}$ 都共面。将 $P$ 与两已知直线分别确定成平面，再联立得到所求直线：
\ans{$\begin{cases}12x-3y-4z-1=0,\\3x+10y-5z-15=0.\end{cases}$}
\end{enumerate}

\Q{9}{过一条直线、平行另一条直线的平面}
题设两线分别是
\[
 l_1:\ \begin{cases}2x+y-z+1=0,\\x-y+2z=0,\end{cases}\quad
 l_2:\ \begin{cases}3x-z+1=0,\\y+2z-2=0.\end{cases}
\]
设过 $l_1$ 的平面族 $\lambda(2x+y-z+1)+\mu(x-y+2z)=0$；由平面平行 $l_2$ 得其法向量与 $l_2$ 方向正交，解得 $\lambda:\mu$。
\ans{$33x+6y+z+13=0$。}

(2) 若求同时与两线共面的、平行 $u=(1,2,1)$ 的直线，可以先分别取过各线且包含 $u$ 的平面，联立为目标直线。\ans{$\begin{cases}x-4y+7z-1=0,\\6x-y-4z+4=0.\end{cases}$}

\Q{10}{共面的直线与平面族的交线}
已知
\[l_1:\begin{cases}x+2y-z+1=0,\\x-4y-z-2=0,\end{cases}\quad
l_2:\begin{cases}x-y+z-2=0,\\4x-2y+1=0.\end{cases}\]
(1) 求过原点、分别与 $l_1,l_2$ 共面的直线。对每个 $l_i$，由直线两方程的线性组合求经过原点的唯一平面，再联立这两个平面：
\ans{$\begin{cases}x-z=0,\\9x-5y+z=0.\end{cases}$}

(2) 过 $l_1$ 且平行 $l_2$ 的平面可求得
\ans{$x-z=0$。}

\Q{11}{一般直线两行列式判别}
把每条直线看作两平面的交线，方向 $u_1=n_1\times n_2$、$u_2=n_3\times n_4$。
\begin{enumerate}[label=(\arabic*)]
\item $l_1\parallel l_2$ 的方向条件是 $u_1\times u_2=0$；将展开式整理，可得到原题写出的两个三阶行列式同为零的判据（前提两条直线都确由非平行平面相交确定）。
\item 若两条直线异面，则四张给定平面没有公共点；四元线性方程组的系数增广矩阵必须满秩，故对应四阶行列式非零。
\item 两线共面意味着没有“独立的四个齐次约束”；按上述逻辑或逆否命题，四阶行列式必须等于零。
\end{enumerate}
\tip{(3) 的行列式为零只是共面的必要条件；想证明充分条件，需要同时考虑平行与相交的几何情况及每条线的非退化假设。}

\Q{12}{两条相交直线与固定方向确定的直线族}
题设
\[l_1:\frac{x-6}{3}=\frac{y}{2}=\frac{z-1}{2},\qquad
 l_2:\frac{x}{3}=\frac{y}{2}=\frac{z+4}{-2}.\]
所有同时与两线相交且平行于 $2x+3y-5=0$ 的直线构成一张直纹二次曲面。设两线上点 $P=P_0+su$、$Q=Q_0+tv$，要求 $(Q-P)\cdot(2,3,0)=0$，由此消去 $s,t$，再令动点 $X=P+w(Q-P)$。把 $s,t,w$ 消去，可得
\ans{$4x^2-9y^2+6x-45y-36z-144=0$，与书后答案一致。}

\Q{13}{三族共面直线决定的二次曲面}
已知 $M_1=(1,0,0),M_2=(-1,0,0),M_3=(2,-1,-2)$，以及
\[u_1=(0,1,1),\ u_2=(0,1,-1),\ u_3=(-3,4,5).\]
(1) 每条已知 $l_i=M_i+t u_i$ 都满足“存在一条与 $l_1,l_2,l_3$ 分别共面的直线过点”的性质：此时第 $i$ 条直线给出的共面线性约束自动为零，只剩另两条关于所求方向的齐次约束，三维方向空间中必有非零解。

(2) 令 $X=(0,1,t)$。过 $X$ 分别与三条已知直线共面的直线，需要其方向 $v$ 同时满足三条齐次线性约束
\[
\det(X-M_i,u_i,v)=0\quad(i=1,2,3).
\]
存在非零 $v$ 的条件是由三行系数组成的矩阵行列式为零，代入 $X=(0,1,t)$ 后可求得\ans{$t=0$。}

(3) 对一般 $X=(x,y,z)$ 同理；消去方向 $v$ 后的行列式展开因子为
\ans{$x^2+y^2-z^2-1=0$，即单叶双曲面。}
